\bookchapter{第十四章：密碼管理器專題}

\chapterintro{在數字時代，密碼就是你的家門鑰匙。但鑰匙太多、太複雜，怎麼辦？密碼管理器就是你的“萬能鑰匙管家”。}

\booksection{一、為什麼你需要密碼管理器？}

現代人平均有幾十甚至上百個網路賬號。如果每個賬號都用相同的密碼，一旦一個網站被黑，所有賬號都危險。如果每個賬號都用不同的複雜密碼，根本記不住。

密碼管理器解決了這個困境：你只需要記住一個主密碼，它幫你記住所有其他密碼。

\booksection{二、密碼管理器的工作原理}

\begin{enumerate}
  \item \textbf{加密儲存}：所有密碼都經過高強度加密儲存，只有你能解密。
  \item \textbf{主密碼}：你設定一個主密碼，用於解鎖密碼庫。主密碼本身不儲存在伺服器上。
  \item \textbf{自動填充}：在網站或App登入介面，密碼管理器自動填充賬號密碼。
  \item \textbf{密碼生成}：可以生成隨機、複雜的密碼，提高安全性。
\end{enumerate}

\booksection{三、加密原理簡介}

瞭解密碼管理器的加密原理，能幫你判斷一個服務商是否值得信任。下面用通俗的語言解釋幾個核心概念。

\booksubsection{AES-256 加密}

AES（高階加密標準）是美國國家標準與技術研究院（NIST）認可的對稱加密演算法，被銀行、軍方、政府廣泛使用。“256”指金鑰長度為 256 位——理論上需要超過宇宙年齡的時間才能暴力破解。主流密碼管理器（1Password、Bitwarden、KeePass 等）都用 AES-256 加密你的密碼庫。

\booksubsection{零知識架構（Zero-Knowledge Architecture）}

“零知識”不是說服務商什麼都不存，而是說\textbf{加密和解密都在你的裝置本地完成，伺服器只儲存已加密的密文}。具體流程：

\begin{enumerate}
  \item 你在客戶端輸入主密碼。
  \item 客戶端用主密碼派生出加密金鑰（這一步在你手機/電腦上完成）。
  \item 只有加密後的密文才上傳到伺服器。
  \item 下次登入時，伺服器把密文下發，客戶端再次用主密碼解密。
\end{enumerate}

關鍵點：\textbf{主密碼本身從不離開你的裝置}，伺服器上也沒有。所以即使服務商被黑，駭客拿到的也只是一堆無法解密的密文——除非他們能猜出你的主密碼。

\begin{tcbinfo}{為什麼"零知識"意味著服務商也看不到你的密碼？}
因為解密金鑰由你的主密碼派生，而主密碼只存在於你腦中、只在你本地裝置上使用。服務商的伺服器上既沒有主密碼，也沒有解密金鑰，只有密文。這就是為什麼 Bitwarden、1Password 等正規服務商都明確宣告“我們也無法看到你的密碼”——不是他們不想看，是技術上做不到。

\end{tcbinfo}

\booksubsection{金鑰派生函式：PBKDF2 與 Argon2}

主密碼不能直接用作加密金鑰——它需要經過“金鑰派生函式”處理。這個函式的作用是\textbf{故意讓計算變慢}，讓駭客暴力破解的成本高到不可行。

\begin{itemize}
  \item \textbf{PBKDF2}：老牌標準，被 iCloud 鑰匙串、Bitwarden、LastPass 等廣泛使用。透過反覆迭代（通常 10 萬次以上）雜湊來拖慢破解。
  \item \textbf{Argon2}：2015 年密碼雜湊競賽冠軍，比 PBKDF2 更現代。除了"算力消耗"還引入了"記憶體消耗"——讓 GPU/ASIC 等專用硬體也難以加速破解。Bitwarden、KeePassXC 等已支援。
\end{itemize}

迭代次數越高越安全，但解鎖速度也會變慢。正規服務商都會公開其使用的派生演算法和迭代次數，便於安全社群審計。

\booksubsection{案例警示：LastPass 2022 年資料洩露}

2022 年 8 月，老牌密碼管理器 LastPass 公告其開發環境被入侵；同年 11 月進一步披露，駭客拿走了使用者的密碼庫備份（加密後的密文）。事件的關鍵教訓：

\begin{itemize}
  \item \textbf{零知識架構守住了底線}：由於密碼庫本身是加密的，駭客拿到密文但解不開——前提是使用者的主密碼足夠強。
  \item \textbf{弱主密碼使用者最危險}：駭客可以離線暴力破解弱主密碼，因此使用弱主密碼或重複主密碼的使用者面臨實質風險。
  \item \textbf{後設資料也洩露}：雖然密碼內容加密，但站點 URL、備註等後設資料可能未加密，會暴露使用者使用了哪些服務。
\end{itemize}

\begin{tcbwarn}{從 LastPass 事件學到的}
第一，主密碼必須長且獨特（建議 16 位以上，包含大小寫、數字、符號）；第二，優先選擇支援 Argon2id、迭代次數高的服務商；第三，對高價值賬號（銀行、郵箱）務必開啟獨立的二次驗證，不要只依賴密碼管理器。

\end{tcbwarn}

\booksection{四、主流密碼管理器推薦}

\booksubsection{系統自帶：}

\begin{itemize}
  \item \textbf{iCloud鑰匙串}：蘋果生態使用者首選。與Safari、App無縫整合，免費，安全。缺點是跨平臺使用不便。
  \item \textbf{Google密碼管理器}：安卓使用者方便。與Chrome、安卓系統整合，免費。
  \item \textbf{華為密碼保險箱}：鴻蒙系統內建，適合華為使用者。
\end{itemize}

\booksubsection{第三方專業工具：}

\begin{itemize}
  \item \textbf{Bitwarden}：開源免費，功能強大，支援所有平臺。強烈推薦給注重隱私的使用者。
  \item \textbf{1Password}：介面美觀，功能豐富，適合家庭和團隊。付費。
  \item \textbf{LastPass}：老牌密碼管理器，免費版功能受限。
  \item \textbf{KeePass}：完全本地儲存，開源免費。適合極客使用者。
\end{itemize}

\booksection{五、密碼強度真實評估}

密碼管理器普遍會顯示“密碼強度”指示條（紅/黃/綠），但這個評估往往只看長度和字元種類，並不全面。真正判斷密碼是否安全，需要考慮“熵”和“是否已洩露”。

\booksubsection{熵（Entropy）：密碼強度的科學度量}

“熵”衡量密碼的不確定性，單位是位元。熵越高，暴力破解所需嘗試次數越多。

\begin{itemize}
  \item \textbf{"123456"}：熵極低，瞬間破解。
  \item \textbf{"P@ssw0rd"}：看似複雜，但出現在所有字典攻擊列表裡，熵依然低。
  \item \textbf{"correct horse battery staple"}：四個隨機單詞，熵約 44 位元，破解難度遠超"P@ssw0rd"。
  \item \textbf{16 位隨機字元（含大小寫+數字+符號）}：熵約 100 位元，目前算力無法暴力破解。
\end{itemize}

結論：\textbf{“長度 + 隨機性”比“複雜符號”更重要}。這也是密碼管理器生成的隨機長密碼真正安全的原因。

\booksubsection{HaveIBeenPwned：查你的密碼是否已洩露}

安全研究員 Troy Hunt 維護的\href{https://haveibeenpwned.com}{haveibeenpwned.com}收錄了歷次大型資料洩露中流出的賬號和密碼。你可以：

\begin{enumerate}
  \item 輸入郵箱，查你的賬號是否在某次洩露中被曝光。
  \item 用 Pwned Passwords 功能查詢某個密碼是否在已知洩露庫中（密碼管理器和瀏覽器都已內建此查詢）。
\end{enumerate}

如果一個密碼出現在洩露庫中，\textbf{無論它看起來多複雜都必須立即更換}——因為字典攻擊會優先嚐試這些已知密碼。

\begin{tcbinfo}{k-Anonymity 查詢原理}
擔心在 HaveIBeenPwned 輸入密碼會再次洩露？不必擔心。現代查詢採用“k-匿名”協議：瀏覽器只把密碼雜湊的前 5 位發給伺服器，伺服器返回所有以這 5 位開頭的雜湊列表，本地再比對。也就是說，伺服器永遠看不到你的完整密碼雜湊。1Password、Bitwarden 的“已洩露密碼檢測”都是這樣工作的。

\end{tcbinfo}

\booksubsection{密碼管理器內建的監控功能}

\begin{itemize}
  \item \textbf{1Password Watchtower}：實時監控弱密碼、重複密碼、已洩露密碼、過期密碼、不使用 2FA 的網站，並以紅/橙/黃分級展示。
  \item \textbf{Bitwarden Vault Health}：免費版即可檢視弱密碼、重複密碼；付費版可檢測已洩露密碼（基於 HaveIBeenPwned）。
  \item \textbf{iCloud 鑰匙串"安全建議"}：iOS 設定 密碼 安全建議，會提示你哪些密碼出現在已知洩露中、哪些重複使用。
  \item \textbf{Google 密碼管理器"密碼檢查工具"}：Chrome 內建，可一鍵檢查已洩露、弱、重複密碼。
\end{itemize}

\begin{tcbpractice}{現在就查：你的密碼是否安全}
\begin{enumerate}
  \item 開啟 haveibeenpwned.com，輸入你的常用郵箱，看看出現在哪些洩露事件中。
  \item iPhone 使用者：設定 密碼 安全建議，逐條處理紅色警告。
  \item Chrome 使用者：設定 密碼 檢查密碼，按提示修改。
  \item 1Password/Bitwarden 使用者：開啟 Watchtower/Health Reports，優先處理"已洩露"和"重複"項。
\end{enumerate}

\end{tcbpractice}

\booksection{六、Passkey：無密碼登入的未來}

2024-2026 年，密碼技術正在經歷一場革命——\textbf{Passkey（通行金鑰）}有望徹底淘汰傳統密碼。Apple、Google、Microsoft 三大平臺聯合推動，FIDO 聯盟和全球資訊網聯盟（W3C）背書，這是數十年來登入方式最大的一次變革。

\booksubsection{什麼是 Passkey？}

Passkey 基於 FIDO2 / WebAuthn 標準，用\textbf{裝置上的金鑰對}替代密碼。原理可以簡化為：

\begin{itemize}
  \item 你的手機生成一對金鑰：\textbf{公鑰}（註冊給網站）+\textbf{私鑰}（永遠存在你裝置的安全晶片中）。
  \item 登入時，網站發一個"挑戰"，你的手機用私鑰簽名，再由網站用公鑰驗證。
  \item 簽名前需要你透過生物識別（Face ID / 指紋）或裝置 PIN 確認身份。
  \item 整個過程中，\textbf{沒有任何密碼在網路中傳輸}。
\end{itemize}

\booksubsection{Passkey vs 傳統密碼}

\begin{center}
\begin{tabularx}{\linewidth}{|X|X|X|}
\hline
\textbf{維度} & \textbf{傳統密碼} & \textbf{Passkey} \tabularnewline
\hline
抗釣魚 & 差，釣魚網站可竊取密碼 & 天然抗釣魚，私鑰只對真實域名簽名 \tabularnewline
伺服器洩露影響 & 密碼雜湊可能被破解 & 伺服器只存公鑰，洩露無影響 \tabularnewline
便捷性 & 需記憶/輸入/管理 & Face ID/指紋一秒登入 \tabularnewline
重放攻擊 & 可能被中間人截獲重放 & 每次簽名一次性，無法重放 \tabularnewline
\hline
\end{tabularx}
\end{center}

\booksubsection{各平臺支援情況}

\begin{itemize}
  \item \textbf{Apple iCloud 鑰匙串 Passkey}：iOS 16+ / macOS Ventura+ 原生支援，Passkey 可在所有登入同一 Apple ID 的裝置間同步，新裝置登入網站時無需重新建立。
  \item \textbf{Google 密碼管理器 Passkey}：安卓 + Chrome 原生支援，Passkey 存在 Google 賬戶中，跨安卓裝置同步；2023 年起 Google 個人賬號已支援用 Passkey 登入。
  \item \textbf{Microsoft}：Windows 11 / Edge 已支援 Passkey；微軟賬號、Outlook、Office 等服務陸續開放 Passkey 登入，並可在 Microsoft Authenticator 中管理。
  \item \textbf{1Password / Bitwarden}：已支援 Passkey 的建立、儲存、跨平臺同步，讓你不受限於單一生態。
\end{itemize}

\begin{tcbinfo}{跨裝置同步是關鍵突破}
早期 Passkey 的最大痛點是“繫結單臺裝置”——換手機就要重新建立。2023 年起 Apple、Google、Microsoft 實現了 Passkey 的雲端同步：iCloud 鑰匙串同步跨 Apple 裝置，Google 密碼管理器同步跨安卓裝置。1Password 等第三方則打通了跨生態同步。這意味著 Passkey 已具備替代密碼的實用性。

\end{tcbinfo}

\booksubsection{如何在手機上建立和使用 Passkey}

\textbf{iOS（iPhone）路徑：}

\begin{enumerate}
  \item 在支援 Passkey 的網站（如 Google、微軟、亞馬遜、eBay、PayPal）賬號安全頁選擇"新增通行金鑰 / Add passkey"。
  \item 系統彈出 Face ID / Touch ID 驗證視窗，透過後 Passkey 自動生成並儲存到 iCloud 鑰匙串。
  \item 下次登入該網站時，輸入使用者名稱後選擇"使用 iPhone 登入"，Face ID 一刷即可。
  \item 管理入口：設定 密碼 找到對應網站，可見 Passkey 條目。
\end{enumerate}

\textbf{安卓路徑：}

\begin{enumerate}
  \item 在支援 Passkey 的網站賬號安全頁選擇"新增通行金鑰"。
  \item Google 密碼管理器彈窗 選擇你的 Google 賬號 螢幕鎖驗證（指紋/圖案/PIN）。
  \item Passkey 自動儲存到 Google 密碼管理器，跨同賬號安卓裝置同步。
  \item 下次登入時，選擇"使用 Passkey"，螢幕鎖一解即可。
  \item 管理入口：安卓設定 密碼、金鑰和自動填充 密碼管理器 google.com。
\end{enumerate}

\booksubsection{Passkey 跨裝置登入}

如果你在 Windows 電腦上登入一個只有 iPhone 上有 Passkey 的網站，可以掃碼授權：

\begin{enumerate}
  \item 網站登入頁選擇"使用 Passkey"。
  \item 選擇"使用其他裝置"。
  \item 電腦顯示二維碼，用 iPhone 掃碼。
  \item iPhone 彈出 Face ID 驗證，確認後網站在電腦上完成登入。
\end{enumerate}

這一過程透過藍芽+網際網路雙重加密通訊，Passkey 私鑰不會離開你的 iPhone。

\begin{tcbpractice}{現在就試：體驗無密碼登入}
\begin{enumerate}
  \item 登入 Google 賬號 安全 通行金鑰 新增通行金鑰，體驗 Face ID/指紋登入。
  \item 登入微軟賬號 安全 高階安全選項 新增通行金鑰。
  \item 在淘寶、京東、eBay、PayPal 等已支援 Passkey 的服務上嚐鮮。
  \item 保留傳統密碼作為後備，逐步遷移高價值賬號到 Passkey。
\end{enumerate}

\end{tcbpractice}

\begin{tcbwarn}{過渡期建議}
Passkey 並未完全取代密碼——許多網站仍處於“Passkey + 密碼”並存階段。建議：高價值賬號優先開啟 Passkey，但\textbf{不要立即刪除傳統密碼}，留作備用；同時繼續使用密碼管理器，因為不是所有網站都支援 Passkey。把 Passkey 看作“密碼管理器的強力補充”而非“立刻替換”。

\end{tcbwarn}

\booksection{七、硬體安全金鑰}

如果你對安全的要求高於普通使用者——比如持有大量加密貨幣、是高管/記者/活動人士、或保護企業核心賬號——硬體安全金鑰是比簡訊驗證碼更可靠的選擇。

\booksubsection{什麼是硬體安全金鑰？}

硬體金鑰是一個物理裝置（形似 U 盤），內部安全晶片儲存你的私鑰，需要物理接觸或 NFC 才能完成驗證。代表產品：

\begin{itemize}
  \item \textbf{YubiKey}（Yubico 公司）：最主流，支援 FIDO2/WebAuthn、U2F、OTP、PIV、OpenPGP，USB-C / USB-A / NFC 等多種介面。5 系列 50-70 美元。
  \item \textbf{Google Titan}：Google 自家出品，USB-C + NFC / 藍芽，約 30-40 美元。
  \item \textbf{OnlyKey}：開源替代品，支援多種協議。
  \item \textbf{國產替代}：如飛天誠信 ePasskey，價格更親民。
\end{itemize}

\booksubsection{為什麼硬體金鑰比簡訊驗證碼安全？}

\begin{itemize}
  \item \textbf{抗釣魚}：和 Passkey 一樣，硬體金鑰只對真實域名響應，釣魚網站無法騙取驗證碼。
  \item \textbf{抗 SIM 劫持}：簡訊驗證碼可被 SIM-swap 攻擊截獲，硬體金鑰無需手機號。
  \item \textbf{抗中間人}：FIDO2 協議每次簽名都包含挑戰值，無法重放。
  \item \textbf{物理隔離}：私鑰永遠在硬體晶片內，惡意軟體即使攻陷你的電腦也讀不到。
\end{itemize}

\booksubsection{適用場景}

\begin{itemize}
  \item 加密貨幣交易所賬號、錢包恢復。
  \item 企業 SSO、遠端辦公 VPN、SSH 登入伺服器。
  \item 郵箱、GitHub、社交媒體等"主賬號"——一旦被攻破可重置所有其他賬號。
  \item 記者、律師、高管等高敏感人群。
\end{itemize}

\booksubsection{如何配置硬體金鑰作為二次驗證}

以 Google 賬號 + YubiKey 為例：

\begin{enumerate}
  \item 購買 YubiKey（建議一次買兩把，一把日常用、一把備用，存放在不同地點）。
  \item 登入 myaccount.google.com 安全 兩步驗證。
  \item 選擇"新增安全金鑰" 把 YubiKey 插入 USB-C 口或貼近手機 NFC 感應區。
  \item 按提示觸控 YubiKey 金屬觸點完成註冊。
  \item 建議註冊第二把備用金鑰。
  \item 在"兩步驗證"設定裡，把硬體金鑰設為首選驗證方式，簡訊驗證碼作為備用。
\end{enumerate}

主流支援的網站：Google、GitHub、Microsoft、Apple ID、Facebook、Twitter、Dropbox、GitLab、AWS、Coinbase、Binance 等。

\begin{tcbinfo}{雙金鑰策略}
硬體金鑰最大的風險是“丟失”。建議：(1) 同時註冊兩把不同的金鑰；(2) 一把隨身攜帶（鑰匙扣），一把存放在家中保險櫃或銀行保管箱；(3) 部分網站允許把“備用恢復碼”列印出來離線儲存。這樣即使主金鑰丟失，也不會被永久鎖在賬號外。

\end{tcbinfo}

\begin{tcbwarn}{購買渠道要謹慎}
務必從官方店鋪或授權代理購買，不要買二手或來源不明的金鑰——理論上預置過的硬體可能已被植入後門。YubiKey 的韌體在出廠後無法更新，但 Yubico 會持續推送安全公告，關注其官網即可。

\end{tcbwarn}

\booksection{八、家庭/團隊密碼共享}

家庭場景下，Netflix、家庭 Wi-Fi、路由器後臺、智慧門鎖密碼需要全家共用；團隊場景下，伺服器、資料庫、社交媒體賬號需要多人協作。\textbf{用微信發密碼、寫在共享文件裡都是高危做法}——一旦聊天記錄或文件許可權洩露，密碼就洩露了。

\booksubsection{1Password 的“共享保險庫”（Vault）}

1Password 是家庭和團隊共享的標杆產品。它的核心機制是\textbf{多個保險庫}：

\begin{itemize}
  \item \textbf{個人保險庫}：只有你能看，如銀行、郵箱等私人賬號。
  \item \textbf{共享保險庫}：家人或團隊成員都能看，如 Netflix、家庭 Wi-Fi、共享信用卡。
  \item \textbf{許可權分級}：團隊版可設定某些成員"只檢視"或"只檢視但不編輯"。
\end{itemize}

家人或團隊成員無需共享主密碼，各自用自己的賬號登入，但可以看到被授權的共享保險庫內容。密碼修改後所有成員實時同步。

\booksubsection{Bitwarden 的“組織”功能}

Bitwarden 的免費版支援個人密碼共享（透過“傳送 Send”功能限時分享單條密碼）；付費版可建立“組織”，把若干條目共享給家庭成員或團隊同事，許可權可細分。Bitwarden 還提供自託管選項，適合對資料主權要求高的企業。

\booksubsection{家庭場景：共享 Netflix、Wi-Fi 等密碼的安全方式}

\begin{itemize}
  \item \textbf{Netflix / 影片網站}：用 1Password 家庭版的"共享保險庫"存登入密碼，家庭成員各自登入即可，無需在群裡發。
  \item \textbf{家庭 Wi-Fi 密碼}：iOS 裝置間可透過 AirDrop 直接共享 Wi-Fi 密碼給來訪的客人，無需念出來；安卓 12+ 也有"附近分享 Wi-Fi"功能。
  \item \textbf{智慧門鎖、智慧家電}：儘量用廠商 App 的"家庭成員"功能為每人單獨建立賬號，而不是共享一個管理員賬號。
  \item \textbf{家用路由器後臺}：管理員密碼存到 1Password 共享保險庫，訪客 Wi-Fi 單獨開一個 SSID。
\end{itemize}

\begin{tcbinfo}{1Password 家庭版的價效比}
1Password Families 約 4.99 美元/月，支援最多 5 位家庭成員，每人都有獨立的私人保險庫 + 4 個共享保險庫。算下來人均不到 1 美元/月，是目前家庭共享體驗最好的方案。Bitwarden 的家庭版 10 美元/年支援 6 人，價效比更高，但介面體驗略遜。

\end{tcbinfo}

\booksubsection{團隊場景：伺服器/資料庫密碼共享}

\begin{itemize}
  \item \textbf{1Password Teams / Business}：可按部門、專案組建立多個保險庫，細粒度許可權控制，所有訪問行為可審計。
  \item \textbf{Bitwarden Teams / Enterprise}：開源方案，可自託管，企業版支援 SSO 整合、目錄同步。
  \item \textbf{職責分離原則}：管理員密碼不應只有一人掌握——設定 2-3 名"應急聯絡人"，確保關鍵人員離職/休假時業務不中斷。
  \item \textbf{離職流程}：員工離職時立即撤銷其密碼管理器訪問許可權，必要時旋轉所有他接觸過的高敏感密碼。
\end{itemize}

\booksubsection{緊急訪問（Emergency Access）：萬一你出事}

想象一個場景：你獨自管理家庭所有重要賬號，但意外住院或不幸離世，家人無法登入銀行、郵箱、保險賬號。\textbf{“緊急訪問”功能}正是為此設計——授權一位信任的家人/朋友，在你無法響應時獲取你的密碼庫。

\textbf{Bitwarden 緊急訪問設定：}

\begin{enumerate}
  \item 登入 Bitwarden Web Vault 設定 緊急訪問。
  \item 新增信任聯絡人的郵箱，邀請對方接受。
  \item 設定"等待期"（建議 7-14 天）：從對方發起請求到你確認的時間視窗。
  \item 如果你在等待期內不響應（拒絕或不操作），到期後對方自動獲得訪問權。
  \item 你隨時可以手動拒絕請求或撤銷緊急訪問許可權。
\end{enumerate}

\textbf{1Password 緊急訪問：}

\begin{enumerate}
  \item 1Password 不直接叫"緊急訪問"，但透過"家庭/團隊組織"功能可實現——把家人加入家庭版後，將部分條目放入"共享保險庫"。
  \item 對個人版使用者，建議把主密碼告知一位高度信任的家人，或寫入密封信封存放於保險櫃。
  \item 1Password 還提供"恢復碼"，開啟雙重認證時務必列印儲存。
\end{enumerate}

\begin{tcbpractice}{現在就設：為家人留一把"備用鑰匙"}
\begin{enumerate}
  \item 選定一位信任的家人作為"緊急聯絡人"。
  \item 在密碼管理器中設定緊急訪問（Bitwarden）或加入家庭共享保險庫（1Password）。
  \item 把主密碼 + 恢復碼寫在紙上，封進信封，存放於家中保險櫃或告知可信家人。
  \item 每半年回顧一次，確認緊急聯絡人是否仍然合適。
\end{enumerate}

\end{tcbpractice}

\booksection{九、密碼管理器使用最佳實踐}

\begin{itemize}
  \item \textbf{設定強主密碼}：主密碼是最後一道防線，務必設定一個長且複雜的密碼，並牢記。
  \item \textbf{啟用雙重認證}：為密碼管理器賬號開啟雙重認證，增加一層保護。
  \item \textbf{定期檢查弱密碼}：大多數密碼管理器有"安全審計"功能，幫你找出弱密碼和重複密碼。
  \item \textbf{不要忘記主密碼}：主密碼無法找回，忘記就意味著丟失所有密碼。建議寫在紙上，放在安全的地方。
  \item \textbf{緊急訪問}：設定緊急聯絡人，萬一你無法訪問，信任的人可以獲取你的密碼。
\end{itemize}

\begin{tcbwarn}{重要提醒}
密碼管理器雖然方便，但也有風險：如果主密碼洩露，所有密碼都暴露了。所以主密碼絕對不能和其他密碼相同，也不能寫在容易被找到的地方。另外，建議選擇有良好安全記錄的服務商。

\end{tcbwarn}

\begin{tcbpractice}{現在就試：開始使用密碼管理器}
\begin{enumerate}
  \item 選擇一個密碼管理器（推薦Bitwarden或系統自帶）。
  \item 註冊賬號，設定一個強主密碼。
  \item 匯入現有密碼（如果有的話）或開始新增新密碼。
  \item 安裝瀏覽器外掛或手機App，體驗自動填充功能。
  \item 檢查安全審計報告，找出需要改進的密碼。
\end{enumerate}

\end{tcbpractice}
